\scanpage{12}
% archive foliation: 11
(en plus des deux diagonales et des deux codiagonales elles-mêmes,
qui font encore quatre autres ensembles à deux éléments, mais groupés par
paires ``indistinguables''\ldots).
On a des bijections canoniques
\[
 \underline{\omega}_S\simeq C,\qquad \underline{\omega}_A\simeq D
\]
provenant du fait qu'on a
\[
 \underline{\omega}_S\simeq\underline{\omega}_{V_S^*}
 \simeq\underline{\omega}_{\{a\}\amalg C}
 \simeq\underline{\omega}_C\simeq C
\]
la convention pour orienter un ensemble à trois éléments $J$ (ici $V_S^*$)
dont un élément $a$ (ici $\underline{a}$) est fixé, en termes du choix d'un
élément \inserted{$x$} du supplémentaire (ici $V_S^*\setminus\{a\}=C$)
étant de prendre l'orientation $(a,x,y)$
\deleted{\unsure{$(\text{où }y\in V_S^*\setminus\{a,x\})$}}
où $\{y\}=J\setminus\{a,x\}$.

On a une relation remarquable
\[
(*)\qquad \underline{\omega}_Q\wedge C\wedge D\simeq\mathbf{1}_{\mathbb{F}_2}
\]
\margin{i.e. une restriction du groupe structural $\mathbb{F}_2^3$ de
$\underline{\omega}_S\times C\times D$
au sous-groupe $\unsure{V_{\underline{3}}|V_2}(\mathbb{F}_2)$ des $x,y,z$ tels que $x+y+z=0$}
qui provient du fait que
$\underline{\omega}_Q,C,D$ sont les images inverses,
dans l'extension $\mathfrak{D}$ de
$\mathfrak{D}/\{1,\underline{a}\}=\mathfrak{D}_{\mathrm{ab}}$
par $\{1,\underline{a}\}$, des trois éléments non nuls de
$\mathfrak{D}_{\mathrm{ab}}$, dont la
\deleted{produit}\inserted{somme} est nulle.
\deleted{Les}\inserted{Pour} des éléments $\varpi,c,d$ de
$\underline{\omega}_Q\times C\times D$ \deleted{\unsure{on définit}} on a
\[
(*)\qquad
 \varpi\wedge c\wedge d=0
 \quad\Longleftrightarrow\quad
 \varpi c d=1\quad\text{dans }\mathfrak{D}
\]
\margin{[Attention, cette relation serait inversée si on faisait une
transposition dans l'ordre des facteurs --- il y a donc
\deleted{\unsure{une}} dans \deleted{\unsure{ce}} le choix de l'iso.
\unsure{ici} une question de convention.}

\margin{Si on passe du carré $Q$ au carré dual $Q'$, le choix de l'iso
\deleted{\unsure{de}}
$\underline{\omega}_{Q'}\wedge C'\wedge D'
\deleted{\unsure{\simeq\mathbf{1}}}
(\simeq\underline{\omega}_Q\wedge D\wedge C)\simeq\mathbf{1}$
est opposé de celui pour $Q$\ldots}
